\documentclass[10pt]{letter}
\usepackage[margin=1in]{geometry}
% Generated by code/analysis.py from data/summary.json. Do not edit by hand.
\newcommand{\secPerUnit}{0.319}
\newcommand{\lamBen}{0.299}
\newcommand{\lamBenEns}{0.296}
\newcommand{\nEns}{47}
\newcommand{\nChaotic}{44}
\newcommand{\nReg}{3}
\newcommand{\lamDir}{0.33}
\newcommand{\lamDirEarly}{0.28}
\newcommand{\lamDirLate}{0.33}
\newcommand{\lamDirOther}{0.329}
\newcommand{\slopeStep}{7.9}
\newcommand{\slopeStepSE}{0.4}
\newcommand{\slopeStepPred}{9.3}
\newcommand{\lamStep}{0.35}
\newcommand{\lamStepSE}{0.02}
\newcommand{\slopePrec}{1.99}
\newcommand{\slopePrecSE}{0.06}
\newcommand{\slopePrecPred}{2.32}
\newcommand{\lamPrec}{0.35}
\newcommand{\lamPrecSE}{0.01}
\newcommand{\gammaRound}{0.6}
\newcommand{\gammaLo}{0.3}
\newcommand{\gammaHi}{0.8}
\newcommand{\lamDirWin}{0.35}
\newcommand{\lamPrecCoarse}{0.31}
\newcommand{\lamPrecFine}{0.34}
\newcommand{\lamExcess}{18}
\newcommand{\TsingleFine}{30}
\newcommand{\collapseSpan}{8}
\newcommand{\collapse}{0.21}
\newcommand{\rkOrder}{4.000}
\newcommand{\ordEuler}{1.002}
\newcommand{\ordMid}{1.999}
\newcommand{\TEulerFine}{31}
\newcommand{\TRKcoarse}{43}
\newcommand{\TkEight}{79}
\newcommand{\TkEleven}{97}
\newcommand{\TBestDouble}{101}
\newcommand{\TBestDoubleSec}{32}
\newcommand{\kBestDouble}{11}
\newcommand{\TBestSingle}{53}
\newcommand{\TBestSingleSec}{17}
\newcommand{\kBestSingle}{7}
\newcommand{\kRatio}{64}
\newcommand{\TdoubleFine}{98}
\newcommand{\TdoubleKten}{92}
\newcommand{\dEexp}{10}
\newcommand{\costFactor}{2000}
\newcommand{\extraBits}{50}
\newcommand{\lamSI}{0.94}
\newcommand{\decadeSI}{10}
\newcommand{\bestRows}{24 (single) & $2^{-7}$ & 53 & 17 \\
32 & $2^{-7}$ & 66 & 21 \\
40 & $2^{-9}$ & 79 & 25 \\
53 (double) & $2^{-11}$ & 101 & 32 \\
64 (extended) & $2^{-15}$ & 124 & 40 \\}

\signature{Arjun Pemmasani\\ on behalf of both authors}
\address{Harvey Mudd College\\ Claremont, CA 91711, USA\\ apemmasani@hmc.edu}
\begin{document}
\begin{letter}{The Editor\\ European Journal of Physics}
\opening{Dear Editor,}

When numerical error grows exponentially at rate $\lambda$, reducing its
effective amplitude by a factor $r$ extends the accuracy horizon by
$\ln r/\lambda$. For RK4 in the truncation-dominated regime, a tenfold smaller
step buys $4\ln10$ Lyapunov times, and a wider floating-point format gives
similarly fixed gains when roundoff dominates. This limitation applies across
chaotic simulations. The double pendulum gives students a familiar system in
which to measure it.

We submit the manuscript ``Step size, floating-point precision and the
predictability horizon of a simulated double pendulum: a student computational
project'' for consideration as a Paper in the Mechanics section of
\emph{European Journal of Physics}.

The project began in our undergraduate numerical analysis course and fits the
journal's interest in original students' computational physics projects. We
vary step size and precision independently, measure truncation and roundoff
separately, and recover the finite-time Lyapunov exponent from the accuracy
horizon's slopes. The project makes an established error-growth argument
accessible through experiments students can run themselves.

\textbf{Intended level.} Advanced undergraduate, in physics or applied
mathematics. The reader needs the fourth-order Runge--Kutta method, the idea of a
floating-point number, and the definition of a Lyapunov exponent.

\textbf{Usefulness to physics education.} Students measure the Lyapunov
exponent through quantities they control in their programs. They learn why a
smaller step can shorten the accuracy horizon and why small energy error does
not certify a trajectory. The manuscript describes classroom use and extensions
to other chaotic systems, including the Lorenz equations and a driven pendulum.

\textbf{Relation to previous work in the journal.} D'Alessio, \emph{Eur.\ J.\
Phys.}\ \textbf{44} (2023) 015002, treats the double pendulum analytically,
numerically and experimentally and computes its Lyapunov spectrum. Our manuscript
is complementary. It takes the dynamics as given and studies the reliability of
the numerical method, which that paper and the earlier pedagogical literature
use as a tool.

The manuscript is below 4000 words including captions by TeXcount, with six figures and two tables. Code, data
and two animations are supplied as supplementary material. The work is original,
has not been published, and is not under consideration elsewhere. Both authors have approved the submission. The authors declare no conflict of interest. The Acknowledgments disclose
the tools used for manuscript revision and code assistance.

\closing{Yours faithfully,}
\end{letter}
\end{document}
